\section{资源账本：参数、激活、优化器状态和带宽}

这门课反复强调 efficiency，不是因为“省钱”这么简单，而是因为语言模型训练真正受制于一张资源账本。
同样的架构放到不同硬件、不同 batch size、不同上下文长度下，主瓶颈可能完全不同。

粗略地说，一个训练系统需要同时关心四类资源：
\begin{itemize}
    \item 参数：模型本身有多大。
    \item 激活：前向传播保留下来的中间状态有多少。
    \item 优化器状态：Adam 之类优化器会额外维护状态。
    \item 带宽与通信：GPU 之间、GPU 与显存之间如何移动数据。
\end{itemize}

\begin{figure}[H]
\centering
\begin{tikzpicture}[every node/.style={font=\small, align=center}]
\node[draw, rounded corners, fill=red!10, minimum width=3.2cm, minimum height=0.9cm] (params) {parameters};
\node[draw, rounded corners, fill=red!18, below left=0.8cm and 1.0cm of params, minimum width=3.2cm, minimum height=0.9cm] (acts) {activations};
\node[draw, rounded corners, fill=red!18, below right=0.8cm and 1.0cm of params, minimum width=3.2cm, minimum height=0.9cm] (opt) {optimizer states};
\node[draw, rounded corners, fill=red!28, below=1.8cm of params, minimum width=3.8cm, minimum height=0.9cm] (bw) {bandwidth / communication};
\draw[-{Latex[length=3mm]}, thick] (params) -- (acts);
\draw[-{Latex[length=3mm]}, thick] (params) -- (opt);
\draw[-{Latex[length=3mm]}, thick] (acts) -- (bw);
\draw[-{Latex[length=3mm]}, thick] (opt) -- (bw);
\end{tikzpicture}
\caption{训练中的资源账本不是单项最优化，而是参数、激活、优化器状态和带宽的共同约束}
\end{figure}

\begin{table}[H]
\centering
\begin{tabular}{p{3.0cm}p{4.2cm}p{5.8cm}}
\toprule
\textbf{资源} & \textbf{常见浪费方式} & \textbf{课程关注的优化方向} \\
\midrule
参数 & 模型过大，无法放入设备 & tensor/pipeline sharding \\
激活 & 上下文过长、batch 太大 & sequence parallelism, checkpointing \\
优化器状态 & Adam 状态膨胀显存占用 & sharding, fused kernels \\
带宽 & GPU 间同步过频繁 & collectives, layout-aware design \\
\bottomrule
\end{tabular}
\caption{资源浪费在不同层的典型形式}
\end{table}

\begin{importantbox}{为什么训练和推理的瓶颈不同}
训练阶段会反复做前向和反向传播，因此更像“算力 + 显存”的联合约束；推理阶段则更容易在 decode 阶段变成“逐 token 生成 + 内存带宽”问题。
\end{importantbox}

\begin{warningbox}{不要被“参数量”误导}
参数量只是容量指标之一。真正决定系统可扩展性的，还有上下文长度、通信开销、优化器状态和推理访问模式。
\end{warningbox}

\subsection{一个简单的预算问题}

假设你有固定预算，能做的事情通常不是“把每一项都拉满”，而是决定把预算花在模型大小、训练 token 还是系统优化上。
这也是为什么这门课每一章都在回答同一个问题：在有限预算下，怎样把边际收益做大？

\begin{figure}[H]
\centering
\begin{tikzpicture}[node distance=1.0cm, every node/.style={font=\small, align=center}]
\node[draw, rounded corners, fill=yellow!12, minimum width=3.2cm, minimum height=0.8cm] (c) {compute};
\node[draw, rounded corners, fill=yellow!20, right=1.2cm of c, minimum width=3.2cm, minimum height=0.8cm] (d) {data};
\node[draw, rounded corners, fill=yellow!28, below=1.1cm of c, minimum width=3.2cm, minimum height=0.8cm] (m) {memory};
\node[draw, rounded corners, fill=yellow!36, right=1.2cm of m, minimum width=3.2cm, minimum height=0.8cm] (b) {bandwidth};
\draw[-{Latex[length=3mm]}, thick] (c) -- (d);
\draw[-{Latex[length=3mm]}, thick] (c) -- (m);
\draw[-{Latex[length=3mm]}, thick] (d) -- (b);
\draw[-{Latex[length=3mm]}, thick] (m) -- (b);
\end{tikzpicture}
\caption{课程中的资源约束图：compute、data、memory、bandwidth 互相牵制}
\end{figure}

